\section{安全性の補題と定理}
\label{sec:systemfxi-safety}

以下は Lean で形式化した命題のステートメントである．証明はここでは記さない．
大域的シグネチャ $\Sigma$ は整形式とする．すなわち，各宣言
$\Sigma(\mathsf{op})=
 \forall\overline{\alpha::\bold{T}}.
 \sigma_{\mathrm{arg}}\to\sigma_{\mathrm{res}}$
の引数型と結果型は，宣言された型変数の下で kind $\bold{T}$ を持つ．
型変数の改名・置換をそれぞれ $\rho_T$，$\tau_T$，
項変数の改名・置換をそれぞれ $\rho_t$，$s_t$ と書く．
置換はすべて，必要に応じて束縛変数を改名する捕獲回避置換とする．
$\mathsf{Val}(e)$ は値判定，$E[e]$ は評価文脈への代入である．
未処理操作を
\[
  \mathsf{Unhandled}(e,\mathsf{op})
  \quad\Longleftrightarrow\quad
  \exists E,\overline{\theta},v.\quad
    e=E[\texttt{perform}~\mathsf{op}
      \bktt{\overline{\theta}}~v]
    \ \land\ \mathsf{Val}(v)
    \ \land\;
    \mathsf{op}\notin\mathbf{handled}(E)
\]
で定義する．証明で用いる $\mathsf{Open}_{\Sigma}(e,\epsilon)$ は，
この条件に加えて $\mathsf{op}\in\epsilon$ と，
$\Sigma$ における操作宣言・型引数の arity の一致を記録する命題である．

\subsection{部分型と置換}

\begin{lemma}[部分型の反射律と推移律]
\label{lem:fxi-subtyping}
整形式な型 $\sigma,\sigma_1,\sigma_2,\sigma_3$ について，
\[
  \Delta\vdash\sigma<:\sigma,
  \qquad
  \frac{\Delta\vdash\sigma_1<:\sigma_2
        \quad\Delta\vdash\sigma_2<:\sigma_3}
       {\Delta\vdash\sigma_1<:\sigma_3}.
\]
\end{lemma}

\begin{lemma}[型変数の改名と置換]
\label{lem:fxi-type-structural}
$\rho_T$ が $\Delta$ の型変数を $\Delta'$ の型変数へ写し，
$\tau_T$ が $\Delta$ の各型変数を $\Delta'$ で整形式な型へ写すとする．
このとき，次の各判定は保存される．
\begin{align*}
  \Delta\vdash\sigma::\bold{T}
    &\Longrightarrow
      \Delta'\vdash\rho_T(\sigma)::\bold{T},
      \quad
      \Delta'\vdash\sigma[\tau_T]::\bold{T},\\
  \Delta\vdash\sigma_1<:\sigma_2
    &\Longrightarrow
      \Delta'\vdash\rho_T(\sigma_1)<:\rho_T(\sigma_2),
      \quad
      \Delta'\vdash\sigma_1[\tau_T]<:\sigma_2[\tau_T].
\end{align*}
\end{lemma}

\begin{lemma}[型付き項の型変数改名・置換]
\label{lem:fxi-typing-type-substitution}
補題 \ref{lem:fxi-type-structural} と同じ条件の下で，
\begin{align*}
  \Delta\mid\Gamma\vdash e:\sigma\mid\epsilon
  &\Longrightarrow
    \Delta'\mid\rho_T(\Gamma)\vdash
      \rho_T(e):\rho_T(\sigma)\mid\epsilon,\\
  \Delta\mid\Gamma\vdash e:\sigma\mid\epsilon
  &\Longrightarrow
    \Delta'\mid\Gamma[\tau_T]\vdash
      e[\tau_T]:\sigma[\tau_T]\mid\epsilon.
\end{align*}
特に，$\Delta\vdash\theta::\bold{T}$ なら
\[
  \Delta,\alpha::\bold{T}\mid\Gamma
    \vdash e:\sigma\mid\epsilon
  \quad\land\quad
  \alpha\notin\ftv{\Gamma}
  \Longrightarrow
  \Delta\mid\Gamma\vdash
    [\alpha\mapsto\theta]e:
    [\alpha\mapsto\theta]\sigma\mid\epsilon.
\]
\end{lemma}

\begin{lemma}[項変数の改名と置換]
\label{lem:fxi-term-substitution}
$\rho_t$ が $\Gamma$ の各変数を同じ型の $\Gamma'$ の変数へ写すとする．
また，$s_t$ が各 $x:\sigma\in\Gamma$ に対して
$\Delta\mid\Gamma'\vdash s_t(x):\sigma\mid\varnothing$
を満たすとする．このとき，
\begin{align*}
  \Delta\mid\Gamma\vdash e:\sigma\mid\epsilon
  &\Longrightarrow
    \Delta\mid\Gamma'\vdash\rho_t(e):\sigma\mid\epsilon,\\
  \Delta\mid\Gamma\vdash e:\sigma\mid\epsilon
  &\Longrightarrow
    \Delta\mid\Gamma'\vdash e[s_t]:\sigma\mid\epsilon.
\end{align*}
特に，
\[
  \frac{\Delta\mid\Gamma,x:\sigma_1
            \vdash e:\sigma_2\mid\epsilon
        \quad
        \Delta\mid\Gamma\vdash v:\sigma_1\mid\varnothing}
       {\Delta\mid\Gamma\vdash
          [x\mapsto v]e:\sigma_2\mid\epsilon}.
\]
\end{lemma}

\begin{proposition}[置換に関する補助等式]
\label{prop:fxi-substitution-equations}
型の恒等改名・恒等置換，改名と置換の合成について，
\begin{align*}
  \operatorname{id}_T(\sigma)&=\sigma,
  &\sigma[\operatorname{id}_T]&=\sigma,\\
  \rho'_T(\rho_T(\sigma))
    &=(\rho'_T\circ\rho_T)(\sigma),
  &(\sigma[\tau_T])[\upsilon_T]
    &=\sigma[\alpha\mapsto(\tau_T(\alpha))[\upsilon_T]],\\
  \rho_T(\sigma[\tau_T])
    &=\sigma[\alpha\mapsto\rho_T(\tau_T(\alpha))],
  &(\rho_T(\sigma))[\tau_T]
    &=\sigma[\alpha\mapsto\tau_T(\rho_T(\alpha))].
\end{align*}
束縛子の下の置換も捕獲を避けて行う．
操作宣言の型を型引数で具体化する置換は，
外側の改名・置換と交換する．
\end{proposition}

\begin{lemma}[宣言型の具体化]
\label{lem:fxi-declaration-instantiation}
$\Delta,\overline{\alpha::\bold{T}}\vdash
\delta::\bold{T}$ かつ
$\Delta\vdash\overline{\theta}::\overline{\bold{T}}$
ならば，
\[
  \Delta\vdash
    [\overline{\alpha}\mapsto\overline{\theta}]
      \delta::\bold{T}.
\]
\end{lemma}

\subsection{型付けの反転と評価文脈}

\begin{lemma}[値の空効果]
\label{lem:fxi-value-pure}
\[
  \mathsf{Val}(v)
  \quad\land\quad
  \Delta\mid\Gamma\vdash v:\sigma\mid\epsilon
  \Longrightarrow
  \Delta\mid\Gamma\vdash v:\sigma\mid\varnothing.
\]
\end{lemma}

\begin{lemma}[閉じた値の canonical forms]
\label{lem:fxi-canonical-forms}
閉じた値が関数型を持つなら項抽象，全称型を持つなら型抽象である：
\begin{align*}
  \mathsf{Val}(v)
  \ \land\
  \varnothing\mid\varnothing\vdash
    v:(\sigma_1\rightarrow~\epsilon~\sigma_2)
      \mid\epsilon_v
  &\Longrightarrow
    \exists\sigma_a,e.\ v=\lambda x:\sigma_a.e,\\
  \mathsf{Val}(v)
  \ \land\
  \varnothing\mid\varnothing\vdash
    v:(\forall\alpha::\bold{T}.\sigma)
      \mid\epsilon_v
  &\Longrightarrow
    \exists e.\ v=\Lambda\alpha::\bold{T}.e.
\end{align*}
\end{lemma}

\begin{lemma}[抽象の型付けの反転]
\label{lem:fxi-abstraction-inversion}
$\Delta\mid\Gamma\vdash
\lambda x:\sigma_a.e:
(\sigma_1\rightarrow~\epsilon~\sigma_2)\mid\epsilon_v$
ならば，ある $\sigma_0,\epsilon_0$ について
\[
  \begin{aligned}
    \Delta\mid\Gamma,x:\sigma_a
      &\vdash e:\sigma_0\mid\epsilon_0,\\
    \Delta\vdash
      (\sigma_a\rightarrow~\epsilon_0~\sigma_0)
      &<: (\sigma_1\rightarrow~\epsilon~\sigma_2).
  \end{aligned}
\]
また，$\Delta\mid\Gamma\vdash
\Lambda\alpha::\bold{T}.e:
(\forall\alpha::\bold{T}.\sigma)\mid\epsilon_v$
ならば，$\alpha\notin\ftv{\Gamma}$ であり，
ある $\sigma_0$ について
\[
  \begin{aligned}
    \Delta,\alpha::\bold{T}\mid\Gamma
      &\vdash e:\sigma_0\mid\varnothing,\\
    \Delta\vdash
      (\forall\alpha::\bold{T}.\sigma_0)
      &<: (\forall\alpha::\bold{T}.\sigma).
  \end{aligned}
\]
\end{lemma}

\begin{lemma}[操作呼出しの型付けの反転]
\label{lem:fxi-perform-inversion}
$\Sigma(\mathsf{op})=
\forall\overline{\alpha::\bold{T}}.
\sigma_{\mathrm{arg}}\to\sigma_{\mathrm{res}}$ とする．
\[
  \Delta\mid\Gamma\vdash
    \texttt{perform}~\mathsf{op}
      \bktt{\overline{\theta}}~v:
      \sigma\mid\epsilon
\]
ならば，$\Delta\vdash\overline{\theta}::\overline{\bold{T}}$ であり，
ある $\epsilon_{\mathrm{arg}}$ について
\[
  \Delta\mid\Gamma\vdash
    v:[\overline{\alpha}\mapsto\overline{\theta}]
      \sigma_{\mathrm{arg}}
      \mid\epsilon_{\mathrm{arg}},
  \qquad
  \Delta\vdash
    [\overline{\alpha}\mapsto\overline{\theta}]
      \sigma_{\mathrm{res}}<:\sigma,
  \qquad
  \mathsf{op}\in\epsilon.
\]
\end{lemma}

\begin{lemma}[評価文脈の改名・差替え・部分式]
\label{lem:fxi-context}
項変数の改名と評価文脈への代入は交換する：
\[
  \rho_t(E[e])=(\rho_t E)[\rho_t(e)].
\]
また，任意の $\sigma_0,\epsilon_0$ について
\[
  \Delta\mid\Gamma\vdash e:\sigma_0\mid\epsilon_0
  \Longrightarrow
  \Delta\mid\Gamma\vdash e':\sigma_0\mid\epsilon_0
\]
が成り立つならば，
\[
  \Delta\mid\Gamma\vdash E[e]:\sigma\mid\epsilon
  \Longrightarrow
  \Delta\mid\Gamma\vdash E[e']:\sigma\mid\epsilon.
\]
さらに，$E[e]$ が型付けできるなら，
ある $\sigma_0,\epsilon_0$ について
$\Delta\mid\Gamma\vdash e:\sigma_0\mid\epsilon_0$ である．
\end{lemma}

\begin{lemma}[多相操作節の具体化]
\label{lem:fxi-clause-instantiation}
$h=\{\texttt{return}~x_r\mapsto e_r,\,
  \mathsf{op}(x;k)\mapsto e_{\mathsf{op}}\}$ とし，
\textsc{THandler} の前提が成り立つとする．
$\Sigma(\mathsf{op})=
\forall\overline{\alpha::\bold{T}}.
\sigma_{\mathrm{arg}}\to\sigma_{\mathrm{res}}$，
$\Delta\vdash\overline{\theta}::\overline{\bold{T}}$ のとき，
\[
  \Delta\mid
  \Gamma,\,
  x:[\overline{\alpha}\mapsto\overline{\theta}]
      \sigma_{\mathrm{arg}},\,
  k:([\overline{\alpha}\mapsto\overline{\theta}]
      \sigma_{\mathrm{res}})
      \rightarrow~\epsilon_{\mathrm{out}}~
      \sigma_{\mathrm{out}}
  \vdash
  [\overline{\alpha}\mapsto\overline{\theta}]
    e_{\mathsf{op}}:
    \sigma_{\mathrm{out}}\mid\epsilon_{\mathrm{out}}.
\]
\end{lemma}

\begin{lemma}[deep 継続の型付け]
\label{lem:fxi-resumption}
$\Delta\mid\Gamma\vdash
\texttt{handle}~
 E[\texttt{perform}~\mathsf{op}\bktt{\overline{\theta}}~v]
 \texttt{with}~h:
 \sigma_{\mathrm{out}}\mid\epsilon_{\mathrm{out}}$
とし，$\Sigma(\mathsf{op})$ は補題
\ref{lem:fxi-clause-instantiation} と同じ宣言であるとする．
\[
  k_E =
    \lambda y:
      [\overline{\alpha}\mapsto\overline{\theta}]
        \sigma_{\mathrm{res}}.
    \texttt{handle}~E[y]~\texttt{with}~h
\]
と置けば，
\[
  \Delta\mid\Gamma\vdash k_E:
  \bigl(
    [\overline{\alpha}\mapsto\overline{\theta}]
      \sigma_{\mathrm{res}}
    \rightarrow~\epsilon_{\mathrm{out}}~
      \sigma_{\mathrm{out}}
  \bigr)\mid\varnothing.
\]
\end{lemma}

\begin{lemma}[操作節の項置換]
\label{lem:fxi-operation-substitution}
$\Delta\mid\Gamma,x:\sigma_{\mathrm{arg}},
 k:\sigma_{\mathrm{res}}\rightarrow~\epsilon~\sigma
 \vdash e_{\mathsf{op}}:\sigma\mid\epsilon$，
$\Delta\mid\Gamma\vdash v:\sigma_{\mathrm{arg}}\mid\varnothing$，
$\Delta\mid\Gamma\vdash k_E:
 (\sigma_{\mathrm{res}}\rightarrow~\epsilon~\sigma)
 \mid\varnothing$ ならば，
\[
  \Delta\mid\Gamma\vdash
    [k\mapsto k_E][x\mapsto v]e_{\mathsf{op}}:
    \sigma\mid\epsilon.
\]
\end{lemma}

\subsection{進行と保存}

\begin{proposition}[基本簡約の型保存]
\label{prop:fxi-redex-preservation}
$\Delta\mid\Gamma\vdash e:\sigma\mid\epsilon$ とする．
EAppAbs，ETAppAbs，EReturn の各規則によって $e\longrightarrow e'$
ならば，$\Delta\mid\Gamma\vdash e':\sigma\mid\epsilon$ である．
また，$\mathsf{op}\notin\mathbf{handled}(E)$ である
EPerform の簡約も同じ型と効果の上界を保存する．
\end{proposition}

\begin{theorem}[Progress]
\label{thm:fxi-progress}
閉じた式について
$\varnothing\mid\varnothing\vdash e:\sigma\mid\epsilon$
ならば，次のいずれかが成り立つ：
\[
  \mathsf{Val}(e)
  \quad\lor\quad
  \exists e'.\ e\longrightarrow e'
  \quad\lor\quad
  \exists E,\mathsf{op},\overline{\theta},v.\quad
    \begin{gathered}
      e=E[\texttt{perform}~\mathsf{op}
          \bktt{\overline{\theta}}~v],\\
      \mathsf{Val}(v),\quad
      \mathsf{op}\notin\mathbf{handled}(E),\quad
      \mathsf{op}\in\epsilon.
    \end{gathered}
\]
最後の場合，操作宣言は $\Sigma$ に存在し，型引数の数はその arity に等しい．
\end{theorem}

\begin{proposition}[空効果の Progress]
\label{prop:fxi-empty-progress}
\[
  \varnothing\mid\varnothing\vdash
    e:\sigma\mid\varnothing
  \Longrightarrow
  \mathsf{Val}(e)
  \ \lor\;
  \exists e'.\ e\longrightarrow e'.
\]
\end{proposition}

\begin{theorem}[Preservation]
\label{thm:fxi-preservation}
\[
  \Delta\mid\Gamma\vdash e:\sigma\mid\epsilon
  \quad\land\quad
  e\longrightarrow e'
  \Longrightarrow
  \Delta\mid\Gamma\vdash e':\sigma\mid\epsilon.
\]
\end{theorem}

\begin{lemma}[多歩の型保存]
\label{lem:fxi-multistep-preservation}
\[
  \Delta\mid\Gamma\vdash e:\sigma\mid\epsilon
  \quad\land\quad
  e\longrightarrow^{*}e'
  \Longrightarrow
  \Delta\mid\Gamma\vdash e':\sigma\mid\epsilon.
\]
\end{lemma}

\begin{theorem}[Effect safety]
\label{thm:fxi-effect-safety}
$\longrightarrow^{*}$ を $\longrightarrow$ の反射推移閉包とする．
$\varnothing\mid\varnothing\vdash e:\sigma\mid\epsilon$
かつ $e\longrightarrow^{*}e'$ ならば，
\[
  \mathsf{Val}(e')
  \quad\lor\quad
  \exists e''.\ e'\longrightarrow e''
  \quad\lor\quad
  \exists E,\mathsf{op},\overline{\theta},v.\quad
    \begin{gathered}
      e'=E[\texttt{perform}~\mathsf{op}
           \bktt{\overline{\theta}}~v],\\
      \mathsf{Val}(v),\quad
      \mathsf{op}\notin\mathbf{handled}(E),\quad
      \mathsf{op}\in\epsilon.
    \end{gathered}
\]
特に，$\epsilon=\varnothing$ ならば，
到達可能な各 $e'$ は値であるか，一歩簡約できる．
\end{theorem}
